% Generated by paper4/manuscript/build.py from paper4/paper/section_limits.md (5d05e774f1a6e2fd): edit the Markdown, not this file.
\section{Limits}
\textbf{Task and data.} One task: three congestion classes 15–120 minutes ahead, from the measured speeds of three cities. Other targets (speed, travel time), longer horizons and other data sources are not tested. The classes come from a congestion factor relative to free flow, with inherited thresholds.

\textbf{Models.} One architecture, Graph WaveNet, with and without its graph, not tuned. The training budget of 1,000 updates was fixed by a pilot and acts as an early stop set on validation, because the curve on unseen nodes has no plateau. No static attribute of the segments is used (in the pilot, the speed limit acted as an identifier of the node and was removed). The training input covers 60–95\% of the observed set, while evaluation uses all of it.

\textbf{Design.} Hidden segments are whole blocks of a grid: about 8 test blocks per configuration in E11, and 8–9 cells at 60\% coverage in E9, hence wide intervals. The 3-point margin is a convention, not a power calculation; no correction for multiple comparisons is applied, as declared in each pre-registration. Machine times come from concurrent processes and are not isolated costs.

\textbf{The analyses without training.} Two ways of using the neighbours were tested, and neither gained over the model without spatial exchange: the calibrated class of a single neighbour at the forecast origin, and a learned aggregation of several neighbours over three hours. Both concern labels in three classes at 15–120 minutes; continuous speeds, incidents and events were not tested. The calendar is the hour of the week; holidays and events are not modelled. The agreement of nearby labels was measured at the same instant, and its value net of the calendar as a forecast: ``shared calendar'' is a reading of the forecasts, not a decomposition of the same-instant agreement.

\textbf{Placement.} Distances are straight-line between segment coordinates, not along the road. The value of proximity was measured on the block placements of E9; that a hidden segment's value depends only on its distance to the nearest sensor, and not on the rule that placed the sensors, is an assumption. Placement rules that use traffic data are not compared.

\textbf{Part A.} The three datasets do not overlap in time or space, so regime and technology cannot be separated (the ring is observed in 2024, the arterials in 2021). The free-flow references use the whole year: Part A compares measurements, it does not forecast. The 19–72 m band was built a posteriori, and the matching share depends on it. The result on lanes rests on 8 sections with three or four lanes, and the rank correlation on 29 sections; no lane threshold is identified. Every statement of Part A is descriptive.

